\setcounter{ex}{0}

% --- CÂU 1 ---
\begin{ex}
	Cho $\cos\alpha = \dfrac{3}{5}$ với $0 < \alpha < \dfrac{\pi}{2}$. Giá trị của $\sin 2\alpha$ bằng
	\choice
	{\True $\dfrac{24}{25}$}
	{$\dfrac{7}{25}$}
	{$-\dfrac{24}{25}$}
	{$\dfrac{12}{25}$}
	\loigiai{
		Vì $0 < \alpha < \dfrac{\pi}{2}$ nên $\sin\alpha > 0$. Ta có:
		$$\sin\alpha = \sqrt{1 - \cos^2\alpha} = \sqrt{1 - \left(\dfrac{3}{5}\right)^2} = \dfrac{4}{5}.$$
		Khi đó:
		$$\sin 2\alpha = 2\sin\alpha\cos\alpha = 2 \cdot \dfrac{4}{5} \cdot \dfrac{3}{5} = \dfrac{24}{25}.$$
	}
\end{ex}

% --- CÂU 2 ---
\begin{ex}
	Cho $\tan\alpha = 2$. Giá trị của biểu thức $\dfrac{3\sin\alpha - \cos\alpha}{\sin\alpha + 2\cos\alpha}$ bằng
	\choice
	{$\dfrac{7}{4}$}
	{\True $\dfrac{5}{4}$}
	{$\dfrac{5}{2}$}
	{$1$}
	\loigiai{
		Do $\tan\alpha = 2$ nên $\cos\alpha \ne 0$. Chia cả tử và mẫu của biểu thức cho $\cos\alpha$, ta được:
		$$\dfrac{3\sin\alpha - \cos\alpha}{\sin\alpha + 2\cos\alpha} = \dfrac{3\tan\alpha - 1}{\tan\alpha + 2} = \dfrac{3 \cdot 2 - 1}{2 + 2} = \dfrac{5}{4}.$$
	}
\end{ex}

% --- CÂU 3 ---
\begin{ex}
	Tập xác định của hàm số $y = \dfrac{1}{\sin x - \cos x}$ là
	\choice
	{$\mathbb{R} \setminus \left\{\dfrac{\pi}{2} + k\pi,\ k \in \mathbb{Z}\right\}$}
	{$\mathbb{R} \setminus \left\{-\dfrac{\pi}{4} + k\pi,\ k \in \mathbb{Z}\right\}$}
	{$\mathbb{R} \setminus \left\{\dfrac{\pi}{4} + k2\pi,\ k \in \mathbb{Z}\right\}$}
	{\True $\mathbb{R} \setminus \left\{\dfrac{\pi}{4} + k\pi,\ k \in \mathbb{Z}\right\}$}
	\loigiai{
		Hàm số xác định khi và chỉ khi:
		$$\sin x - \cos x \ne 0 \iff \sqrt{2}\sin\left(x - \dfrac{\pi}{4}\right) \ne 0 \iff x - \dfrac{\pi}{4} \ne k\pi \iff x \ne \dfrac{\pi}{4} + k\pi\ (k \in \mathbb{Z}).$$
		Vậy tập xác định của hàm số là $\mathscr{D} = \mathbb{R} \setminus \left\{\dfrac{\pi}{4} + k\pi,\ k \in \mathbb{Z}\right\}$.
	}
\end{ex}

% --- CÂU 4 ---
\begin{ex}
	Hàm số nào dưới đây là hàm số chẵn?
	\choice
	{$y = \sin x + \cos x$}
	{\True $y = x\sin x$}
	{$y = x\cos x$}
	{$y = \tan x + x$}
	\loigiai{
		Xét hàm số $f(x) = x\sin x$ có tập xác định $\mathscr{D} = \mathbb{R}$ là tập đối xứng.\\
		Với mọi $x \in \mathbb{R}$, ta có $f(-x) = (-x)\sin(-x) = (-x)(-\sin x) = x\sin x = f(x)$.\\
		Vậy $y = x\sin x$ là hàm số chẵn.
	}
\end{ex}

% --- CÂU 5 ---
\begin{ex}
	Nghiệm của phương trình $2\sin\left(2x - \dfrac{\pi}{3}\right) = \sqrt{3}$ là
	\choice
	{$x = \dfrac{\pi}{3} + k\pi$ hoặc $x = \dfrac{2\pi}{3} + k\pi\ (k \in \mathbb{Z})$}
	{$x = \dfrac{\pi}{3} + k2\pi$ hoặc $x = \dfrac{2\pi}{3} + k2\pi\ (k \in \mathbb{Z})$}
	{\True $x = \dfrac{\pi}{3} + k\pi$ hoặc $x = \dfrac{\pi}{2} + k\pi\ (k \in \mathbb{Z})$}
	{$x = \dfrac{\pi}{6} + k\pi$ hoặc $x = \dfrac{\pi}{2} + k\pi\ (k \in \mathbb{Z})$}
	\loigiai{
		Ta có:
		$$2\sin\left(2x - \dfrac{\pi}{3}\right) = \sqrt{3} \iff \sin\left(2x - \dfrac{\pi}{3}\right) = \dfrac{\sqrt{3}}{2} \iff \hoac{2x - \dfrac{\pi}{3} &= \dfrac{\pi}{3} + k2\pi \\ 2x - \dfrac{\pi}{3} &= \dfrac{2\pi}{3} + k2\pi} \iff \hoac{x &= \dfrac{\pi}{3} + k\pi \\ x &= \dfrac{\pi}{2} + k\pi}\ (k \in \mathbb{Z}).$$
	}
\end{ex}

% --- CÂU 6 ---
\begin{ex}
	Số nghiệm của phương trình $2\cos 2x - 1 = 0$ trên đoạn $[0; 2\pi]$ là
	\choice
	{$2$}
	{$3$}
	{\True $4$}
	{$6$}
	\loigiai{
		Ta có:
		$$2\cos 2x - 1 = 0 \iff \cos 2x = \dfrac{1}{2} \iff 2x = \pm\dfrac{\pi}{3} + k2\pi \iff x = \pm\dfrac{\pi}{6} + k\pi\ (k \in \mathbb{Z}).$$
		Xét các nghiệm trên đoạn $[0; 2\pi]$:
		\begin{itemize}
			\item Với $x = \dfrac{\pi}{6} + k\pi \in [0; 2\pi] \iff -\dfrac{1}{6} \le k \le \dfrac{11}{6}$. Vì $k \in \mathbb{Z}$ nên $k \in \{0; 1\}$ (có $2$ nghiệm).
			\item Với $x = -\dfrac{\pi}{6} + k\pi \in [0; 2\pi] \iff \dfrac{1}{6} \le k \le \dfrac{13}{6}$. Vì $k \in \mathbb{Z}$ nên $k \in \{1; 2\}$ (có $2$ nghiệm).
		\end{itemize}
		Vậy phương trình có $4$ nghiệm trên đoạn $[0; 2\pi]$.
	}
\end{ex}

% --- CÂU 7 ---
\begin{ex}
	Giá trị nhỏ nhất của hàm số $y = 4\sin\left(x + \dfrac{\pi}{6}\right) - 1$ bằng
	\choice
	{$3$}
	{$-4$}
	{\True $-5$}
	{$-3$}
	\loigiai{
		Vì $-1 \le \sin\left(x + \dfrac{\pi}{6}\right) \le 1, \forall x \in \mathbb{R}$ nên:
		$$-5 \le 4\sin\left(x + \dfrac{\pi}{6}\right) - 1 \le 3.$$
		Vậy giá trị nhỏ nhất của hàm số bằng $-5$ khi $\sin\left(x + \dfrac{\pi}{6}\right) = -1$.
	}
\end{ex}

% --- CÂU 8 ---
\begin{ex}
	Cho cấp số cộng $(u_n)$ có $u_2 = 5$ và $u_5 = 14$. Số hạng tổng quát của cấp số cộng là
	\choice
	{$u_n = 3n + 2$}
	{$u_n = 3n - 2$}
	{$u_n = 2n + 1$}
	{\True $u_n = 3n - 1$}
	\loigiai{
		Gọi $d$ là công sai của cấp số cộng. Theo giả thiết, ta có hệ phương trình:
		$$\heva{u_2 &= u_1 + d = 5 \\ u_5 &= u_1 + 4d = 14} \iff \heva{u_1 &= 2 \\ d &= 3.}$$
		Số hạng tổng quát của cấp số cộng là:
		$$u_n = u_1 + (n - 1)d = 2 + (n - 1) \cdot 3 = 3n - 1.$$
	}
\end{ex}

% --- CÂU 9 ---
\begin{ex}
	Cho cấp số cộng $(u_n)$ có $u_1 = -3$ và công sai $d = 4$. Tổng $12$ số hạng đầu bằng
	\choice
	{\True $228$}
	{$216$}
	{$240$}
	{$204$}
	\loigiai{
		Tổng $12$ số hạng đầu của cấp số cộng là:
		$$S_{12} = 12u_1 + \dfrac{12 \cdot 11}{2}d = 12 \cdot (-3) + 66 \cdot 4 = -36 + 264 = 228.$$
	}
\end{ex}

% --- CÂU 10 ---
\begin{ex}
	Cho hình chóp $S.ABCD$ có đáy $ABCD$ là hình bình hành tâm $O$, $M$ là trung điểm của $SC$ (hình vẽ dùng chung cho câu 10, 11, 12). Giao tuyến của hai mặt phẳng $(SAD)$ và $(SBC)$ là đường thẳng đi qua $S$ và
	\begin{center}
		\begin{tikzpicture}[scale=0.85, font=\footnotesize, line join=round, line cap=round]
			\coordinate (A) at (0,0);
			\coordinate (B) at (3.2,0);
			\coordinate (D) at (1.3,1.1);
			\coordinate (C) at (4.5,1.1);
			\coordinate (O) at ($(A)!0.5!(C)$);
			\coordinate (S) at (2.2,3.4);
			\coordinate (M) at ($(S)!0.5!(C)$);
			\draw[thick] (S)--(A) (S)--(B) (S)--(C) (A)--(B)--(C);
			\draw[dashed] (S)--(D) (A)--(D)--(C) (A)--(C) (B)--(D);
			\foreach \p/\pos in {S/above, A/below left, B/below, C/right, D/above left, O/below, M/above right}
				\node[\pos] at (\p) {$\p$};
			\foreach \p in {S,A,B,C,D,O,M}
				\fill (\p) circle (1.2pt);
		\end{tikzpicture}
	\end{center}
	\choice
	{song song với $AB$}
	{song song với $AC$}
	{đi qua $O$}
	{\True song song với $AD$}
	\loigiai{
		Hai mặt phẳng $(SAD)$ và $(SBC)$ có điểm chung là $S$ và lần lượt chứa hai đường thẳng song song $AD, BC$ ($AD \parallel BC$).\\
		Do đó giao tuyến của $(SAD)$ và $(SBC)$ là đường thẳng đi qua $S$ và song song với $AD$ (và $BC$).
	}
\end{ex}

% --- CÂU 11 ---
\begin{ex}
	Gọi $F$ là giao điểm của đường thẳng $SD$ và mặt phẳng $(ABM)$. Khẳng định nào sau đây đúng?
	\choice
	{\True $F$ là trung điểm của $SD$}
	{$FD = 2SF$}
	{$SF = 2FD$}
	{$SF = 3FD$}
	\loigiai{
		Hai mặt phẳng $(ABM)$ và $(SCD)$ có điểm chung là $M$ và lần lượt chứa hai đường thẳng song song $AB, CD$ ($AB \parallel CD$).\\
		Suy ra giao tuyến của $(ABM)$ và $(SCD)$ là đường thẳng qua $M$ song song với $CD$, cắt $SD$ tại $F$.\\
		Trong tam giác $SCD$, vì $M$ là trung điểm của $SC$ và $MF \parallel CD\ (F \in SD)$ nên $F$ là trung điểm của $SD$.
	}
\end{ex}

% --- CÂU 12 ---
\begin{ex}
	Khẳng định nào sau đây đúng?
	\choice
	{\True $OM \parallel (SAB)$}
	{$OM \parallel (ABCD)$}
	{$OM \parallel (SCD)$}
	{$OM \parallel (SBC)$}
	\loigiai{
		Xét tam giác $SAC$ có $O$ là trung điểm của $AC$ và $M$ là trung điểm của $SC$.\\
		Suy ra $OM$ là đường trung bình của tam giác $SAC \implies OM \parallel SA$.\\
		Mà $SA \subset (SAB)$ và $OM \not\subset (SAB)$ nên $OM \parallel (SAB)$.
	}
\end{ex}
